\subsection*{一步一步建立直观}

\paragraph{1. 从中心开始看两个根的对称性}

对一元二次式，配方把一次项吸收到平移后的平方中，两个实根便围绕同一个中心对称。右图的对称矩阵提供一个相关例子：其特征值满足二次特征方程，两个根的中点是迹的一半。这里先读矩阵，再把它当作检验配方结果的具体对象。



\paragraph{2. 中心加上偏移得到一个根}

当前矩阵的对角元相等，特征方向之一沿两坐标同号的对角线，放大倍率为三。把特征方程写成围绕中心二的平方关系，就得到向右偏移一的根。一般配方的中心与偏移应从系数计算，而不是凭图上看起来对称直接猜出。



\paragraph{3. 另一个根来自相反的偏移}

转到另一条特征方向，放大倍率为一，恰好在中心二的另一侧。这个例子有两个实根，因为实对称矩阵的特征值为实数；普通二次方程还可能出现负判别式。图示说明配方的对称结构，不能把本例的实根性质推广到任意系数。


